\section{Configuración del rack}

Una vez realizados los procesos de reseteo de los dispositivos de red, se procedió con la configuración lógica del rack de comunicaciones. Esta etapa es fundamental, ya que define el comportamiento de la infraestructura de red, establece la segmentación mediante VLANs y permite la interconexión efectiva entre las distintas redes diseñadas.

La configuración del rack incluye la parametrización del router, el switch y el Access Point, asegurando que cada dispositivo cumpla con su función específica dentro de la topología planteada por el Grupo 6.

\subsection{Configuración del Router}

El router Cisco fue configurado para cumplir la función de enrutamiento inter-VLAN, actuando como puerta de enlace predeterminada para cada una de las redes lógicas implementadas. Para ello, se utilizó el esquema conocido como \textit{Router-on-a-Stick}, el cual permite manejar múltiples VLANs a través de una única interfaz física mediante subinterfaces y encapsulación IEEE 802.1Q.

Adicionalmente, el router fue configurado para proporcionar servicios de asignación dinámica de direcciones IP mediante DHCP, así como parámetros básicos de seguridad para el acceso administrativo.

\subsubsection{Configuración básica e interfaces}

\begin{lstlisting}[caption={Configuración de subinterfaces y direccionamiento IP en el Router}]
enable
configure terminal
hostname R2
interface fa0/0
no shutdown
exit
interface fa0/0.10
encapsulation dot1Q 10
ip address 192.168.61.1 255.255.255.0
exit
interface fa0/0.20
encapsulation dot1Q 20
ip address 192.168.62.1 255.255.255.0
exit
interface fa0/0.30
encapsulation dot1Q 30
ip address 192.168.63.1 255.255.255.0
exit
interface fa0/0.40
encapsulation dot1Q 40
ip address 192.168.64.1 255.255.255.0
exit
\end{lstlisting}

Cada subinterface fue asociada a una VLAN específica y configurada con la dirección IP correspondiente, permitiendo que el router actúe como puerta de enlace predeterminada para los dispositivos conectados en cada red.

\subsubsection{Configuración de acceso y seguridad}

\begin{lstlisting}[caption={Configuración de seguridad y acceso remoto en el Router}]
line vty 0 15
password cisco
login
enable secret cisco123
service password-encryption
exit
\end{lstlisting}

Esta configuración permite el acceso remoto al router mediante líneas VTY, asegurando que las credenciales se encuentren cifradas y protegidas frente a accesos no autorizados.

\subsubsection{Configuración del servicio DHCP}

\begin{lstlisting}[caption={Configuración de pools DHCP para cada VLAN}]
configure terminal
ip dhcp pool VLAN10
network 192.168.61.0 255.255.255.0
default-router 192.168.61.1
dns-server 8.8.8.8
exit
ip dhcp pool VLAN20
network 192.168.62.0 255.255.255.0
default-router 192.168.62.1
dns-server 8.8.8.8
exit
ip dhcp pool VLAN30
network 192.168.63.0 255.255.255.0
default-router 192.168.63.1
dns-server 8.8.8.8
exit
ip dhcp pool VLAN40
network 192.168.64.0 255.255.255.0
default-router 192.168.64.1
dns-server 8.8.8.8
exit
exit
copy running-config startup-config
\end{lstlisting}

Cada \textit{pool} DHCP fue definido para su respectiva VLAN, permitiendo la asignación automática de direcciones IP, puerta de enlace y servidor DNS a los equipos finales.

\subsection{Configuración del Switch}

El switch Cisco fue configurado como dispositivo de capa 2 encargado de la segmentación de la red mediante VLANs. Cada puerto fue asignado según el tipo de dispositivo conectado, utilizando puertos de acceso para equipos finales y puertos troncales para la comunicación con el router.

\subsubsection{Creación de VLANs}

\begin{lstlisting}[caption={Creación y nombramiento de VLANs en el Switch}]
enable
configure terminal
hostname S1
vlan 10
name KALI_ACCESS
exit
vlan 20
name SERVER
exit
vlan 30
name PCG6
exit
vlan 40
name PRUEBAS
exit
\end{lstlisting}

La creación de VLANs permite separar lógicamente el tráfico de red, mejorando tanto la seguridad como el rendimiento de la infraestructura.

\subsubsection{Asignación de puertos de acceso}

\begin{lstlisting}[caption={Configuración de puertos de acceso en el Switch}]
interface fa0/1
switchport mode access
switchport access vlan 10
exit
interface fa0/2
switchport mode access
switchport access vlan 20
exit
interface fa0/3
switchport mode access
switchport access vlan 30
exit
interface fa0/4
switchport mode access
switchport access vlan 40
exit
\end{lstlisting}

Cada puerto fue configurado en modo acceso y asignado a la VLAN correspondiente, conectando de forma directa los equipos finales y el Access Point.

\subsubsection{Configuración de enlaces troncales}

\begin{lstlisting}[caption={Configuración de puertos troncales en el Switch}]
interface fa0/24
switchport mode trunk
switchport trunk allowed vlan 10,20,30,40
no shutdown
exit
interface fa0/23
switchport mode trunk
switchport trunk allowed vlan 10,20,30,40
no shutdown
exit
\end{lstlisting}

Los puertos troncales permiten el transporte de tráfico etiquetado entre el switch y el router, facilitando el enrutamiento inter-VLAN.

\subsubsection{Configuración de acceso administrativo}

\begin{lstlisting}[caption={Configuración de seguridad básica en el Switch}]
line vty 0 15
password cisco
login
exit
enable secret cisco123
service password-encryption
exit
copy running-config startup-config
\end{lstlisting}

Esta configuración asegura el acceso administrativo al switch y protege las credenciales almacenadas.

\subsection{Configuración del Access Point}

El Access Point fue integrado a la red como parte de la VLAN destinada a Kali Linux y pruebas inalámbricas. Su función principal es proporcionar conectividad WiFi y permitir la ejecución de auditorías de seguridad sobre redes inalámbricas en un entorno controlado.

El Access Point fue configurado para operar dentro del segmento correspondiente, obteniendo su dirección IP de forma dinámica mediante DHCP y permitiendo la asociación de dispositivos autorizados para las pruebas académicas.

\subsubsection{Configuración lógica y de seguridad del Access Point}

El Access Point Cisco modelo \textbf{WAP4410N} fue configurado mediante la carga de un archivo de configuración, lo cual permitió establecer de manera consistente los parámetros de red, seguridad inalámbrica y administración del dispositivo. Este enfoque garantiza reproducibilidad y reduce errores derivados de configuraciones manuales.

La configuración aplicada responde a los requerimientos del entorno académico planteado, integrando el Access Point dentro de la infraestructura VLAN previamente definida y habilitando una red inalámbrica segura para la ejecución de pruebas controladas.

\paragraph{Configuración de red y direccionamiento}

El Access Point fue configurado para obtener su direccionamiento IP de forma automática mediante DHCP, permitiendo su integración directa con el router configurado previamente como servidor de direcciones.

\begin{lstlisting}[caption={Configuración básica de red del Access Point}]
[BASIC]
Host_Name=wap1627c4
Device_Name=WAP4410N
IP_settings=1
IPv6=0
\end{lstlisting}

Esta configuración permite una administración flexible del dispositivo, evitando conflictos de direccionamiento y facilitando su despliegue dentro del rack de comunicaciones.

\paragraph{Configuración inalámbrica básica}

El Access Point fue configurado para operar en modo mixto \textit{802.11 b/g/n}, asegurando compatibilidad con distintos dispositivos cliente. Se definió un único SSID activo correspondiente al grupo de trabajo.

\begin{lstlisting}[caption={Configuración inalámbrica básica}]
[Wireless_Basic]
Network_mode=7
Channel=6
SSID1=GRUPO6
SSID1_broadcast=1
\end{lstlisting}

El uso de un canal fijo permite controlar interferencias durante las pruebas de auditoría WiFi, facilitando la captura de tráfico y el análisis posterior.

\paragraph{Configuración de seguridad inalámbrica}

Para garantizar la protección de la red inalámbrica, se implementó el estándar \textbf{WPA2-Personal} con cifrado \textbf{AES}, considerado una buena práctica en entornos modernos. La autenticación se realiza mediante una clave precompartida (PSK).

\begin{lstlisting}[caption={Configuración de seguridad inalámbrica WPA2}]
[Wireless_security_1]
security_mode=3
authentication_type=1
encryption=2
PSK_key=grupo6_cripto
Key_renew=3600
\end{lstlisting}

Esta configuración permite evaluar la robustez de la contraseña durante los ataques documentados en capítulos posteriores, manteniendo un equilibrio entre seguridad y viabilidad de pruebas académicas.

\paragraph{Aislamiento y control de acceso}

Se habilitó el aislamiento entre SSID y se mantuvo deshabilitado el control de acceso por direcciones MAC, con el objetivo de no interferir en las pruebas de auditoría y permitir la asociación controlada de los dispositivos del laboratorio.

\begin{lstlisting}[caption={Parámetros de aislamiento y control de acceso}]
isolation_between_SSID=1
isolation_within_SSID=0
Connection_Control=0
\end{lstlisting}

\paragraph{Parámetros avanzados y operación del dispositivo}

El Access Point fue configurado para operar exclusivamente en modo \textit{Access Point}, descartando funcionalidades adicionales como WDS o modo monitor, ya que la captura de tráfico se realizó mediante una antena externa dedicada.

\begin{lstlisting}[caption={Modo de operación del Access Point}]
[AP_Mode]
AP_mode=0
\end{lstlisting}

Asimismo, se establecieron parámetros avanzados de transmisión, como ancho de canal de 20 MHz y valores estándar de \textit{beacon} y \textit{DTIM}, garantizando estabilidad y rendimiento durante las pruebas.

\paragraph{Configuración administrativa}

Finalmente, se configuraron las credenciales administrativas básicas del dispositivo, restringiendo el acceso a personal autorizado del grupo de trabajo.

\begin{lstlisting}[caption={Configuración administrativa del Access Point}]
[Management]
username=admin
AP_password=admin
HTTPS_Access=0
Secure_Shell=0
\end{lstlisting}

\paragraph{Relación con las pruebas de seguridad}

La configuración del Access Point constituye un elemento clave para el desarrollo del tercer ataque documentado en este trabajo, ya que define el entorno inalámbrico objetivo sobre el cual se realizaron las pruebas de auditoría WiFi, captura de handshakes y análisis de seguridad de la clave precompartida.

\subsection{Importancia de la configuración del rack}

La correcta configuración del rack garantiza el funcionamiento adecuado de la infraestructura de red y permite la ejecución efectiva de las actividades de seguridad planteadas en el proyecto. La segmentación mediante VLANs, el enrutamiento inter-VLAN y la asignación dinámica de direcciones IP constituyen elementos clave para simular un entorno realista y seguro.

Esta etapa sienta las bases para el desarrollo de los ataques documentados en los capítulos posteriores, asegurando que la red se encuentre correctamente estructurada y operativa.
